% chapters/06/section-06.tex
\section{Yang Belum Ada pada Fisika Ini}
\label{sec:belum-ada-fisika}

Benda-benda pada bab ini sudah jatuh, melompat, dan memantul, tetapi masih
terlalu sederhana untuk sebuah permainan yang utuh. Bagian penutup ini
menjelaskan batas kemampuannya.

\subsection{Tabrakan yang Baru Diketahui Sebagian}

Program pada bab ini hanya mengenal dua jenis batas: lantai sebagai satu garis
mendatar, dan kedua dinding sebagai dua garis tegak. Pemeriksaannya pun
dilakukan dengan membandingkan satu nilai, bukan dengan memeriksa bentuk.

Permainan yang sesungguhnya memiliki rintangan di tengah ruang permainan,
benda yang bertumpuk, dan benda yang bergerak saling mendekat. Pemeriksaan
semacam itu menuntut cara yang lebih umum: membandingkan persegi pembatas dua
benda, menentukan dari sisi mana tabrakan terjadi, lalu menyelesaikannya per
sumbu. Seluruhnya dibahas pada Bab~VII.

Pekerjaan pada bab ini bukan pekerjaan yang sia-sia. Urutan penanganan yang
sudah Anda pelajari, yaitu perbaiki posisi lebih dahulu baru ubah kecepatan,
tetap berlaku pada pemeriksaan tabrakan yang lebih umum. Begitu pula kebiasaan
memeriksa dengan tanda ketidaksamaan, bukan kesamaan.

\subsection{Sifat Benda yang Belum Dimodelkan}

Tiga sifat berikut belum ada pada model ini, dan masing-masing menuntut
tambahan tersendiri.

\begin{itemize}
  \item \textbf{Gesekan} (\istilah{friction}). Benda yang bergeser di atas
        lantai seharusnya melambat dan akhirnya berhenti. Pada model ini
        kecepatan mendatarnya hanya berubah karena tombol, sehingga pemain
        dapat meluncur selamanya.
  \item \textbf{Hambatan udara.} Benda yang jatuh sangat lama seharusnya
        mencapai kecepatan akhir dan tidak lagi bertambah cepat. Pada model
        ini kecepatannya bertambah tanpa batas, dan itulah sebabnya penambahan
        batas kecepatan diperlukan untuk mencegah benda menembus lantai.
  \item \textbf{Permukaan yang miring dan bertingkat.} Lantai pada model ini
        hanya berupa satu garis mendatar. Permukaan miring menuntut
        perhitungan arah normal, sedangkan permukaan bertingkat biasanya
        diselesaikan dengan membagi perpindahan menjadi beberapa langkah
        kecil.
\end{itemize}

\subsection{Perhitungan yang Diserahkan kepada Mesin}

Seperti halnya seluruh Blok~I, bab ini memperlihatkan apa yang sesungguhnya
dikerjakan oleh mesin permainan. Pada Unity nanti, gravitasi, pantulan, dan
tabrakan dihitung oleh \istilah{engine} dengan cara yang jauh lebih cermat,
termasuk gesekan, kecepatan akhir, dan pemisahan per sumbu.

Yang tidak berubah adalah cara berpikirnya: benda memiliki posisi dan
kecepatan, percepatan mengubah kecepatan, kecepatan mengubah posisi, dan
tabrakan menyelesaikan posisi lebih dahulu sebelum mengubah kecepatan.
Pengetahuan itu yang membuat Anda dapat memilih pengaturan yang tepat pada
mesin, dan mengenali gejala ketika pengaturannya kurang sesuai.

\subsection{Ringkasan Perkembangan}

\begin{table}[htbp]
\centering
\caption{Bagian yang ditambahkan pada fisika ini pada bab-bab berikutnya}
\label{tab:perkembangan-fisika}
{\setlength{\tabcolsep}{4pt}\footnotesize%
\begin{tabularx}{\textwidth}{@{}C{1.6cm} L{5.0cm} Y@{}}
\toprule
\headrow \textbf{Bab} & \textbf{Yang ditambahkan} & \textbf{Akibatnya} \\
\midrule
Bab~VII & Deteksi tabrakan antar-objek, \code{Sprite} dan \code{Group} &
Benda dapat bertabrakan dengan rintangan dan dengan benda lain, bukan hanya
dengan tepi jendela. \\
\addlinespace
Bab~VIII--XI & Pengendalian NPC &
Kecepatan tidak lagi berasal dari tombol atau bernilai tetap, melainkan dari
keputusan yang diambil program sendiri. \\
\addlinespace
Bab~XIII & Fisika mesin pada Unity &
Perhitungan percepatan, tabrakan, dan gesekan diserahkan kepada mesin,
sedangkan kode Anda mengatur pengaturan dan menanggapi peristiwanya. \\
\bottomrule
\end{tabularx}}
\end{table}

Dengan gambaran tersebut, tampak bahwa bab ini melengkapi tingkat kedua pada
tahap pembaruan. Sesudah bab ini, kerangka perhitungan geraknya lengkap:
percepatan mengubah kecepatan, kecepatan mengubah posisi, dan tabrakan
menyelesaikan keduanya. Bab-bab berikutnya tidak menambah tingkat baru,
melainkan memakai kerangka itu untuk keperluan yang berbeda.
